\input{../header}

\begin{document}
%


%\onehalfspacing
\allowdisplaybreaks
%##################################################################
\section{PS\#12 - Taylor polynomials - \red{Answer key} }

\everymath{\displaystyle}
\begin{enumerate}[leftmargin=0pt]

\item (\href{https://faculty.gvsu.edu/boelkinm/Home/ACS/sec-8a-2-taylor-poly.html#ez-8a-2-0}{AC 8a.2 Exercise 5}) Suppose we know the following information about a function $f$:
\begin{equation*}
f(0) = 2, f'(0) = -3, f''(0) = -1, f'''(0) = 0, f^{(4)}(0) = -3\text{.}
\end{equation*}
\begin{enumerate}
    \item Determine $T_4(x)$, the degree 4 Taylor polynomial of $f$ that is centered at $a=0$.    \begin{red}
    \[T_4(x)=2 - 3x - \frac12 x^2 - \frac18 x^4\]
    \end{red}
    \item Use $T_4(x)$ to estimate $f(0.5)$.
    \begin{red}
    $f(0.5) \approx T_4(0.5) = 0.3671875.$
    \end{red}
    
    \item State each of the Taylor polynomials $T_3(x)$, $T_2(x)$, and $T_1(x)$.
    \begin{red}
    \begin{align*}
        T_4(x) &= 2 - 3x - \frac12 x^2 - \frac18 x^4 \\
        T_3(x) &= 2 - 3x - \frac12 x^2 (+0x^3) \\
        T_2(x) &= 2 - 3x - \frac12 x^2 \\
        T_1(x) &= 2 - 3x
    \end{align*}
    \end{red}
\end{enumerate}

\item (\href{https://faculty.gvsu.edu/boelkinm/Home/ACS/sec-8a-2-taylor-poly.html#ez-8a-2-2}{AC 8a.2 Exercise 7}) 
\begin{enumerate}
    \item Determine the degree 6 Taylor polynomial approximation of $\sin(x)$ centered at $a=\tfrac\pi2$.
    
    \begin{red}
    $\sin(x) \approx 1 
    - \frac{1}{2!} \left(x-\frac\pi2\right)^2
    + \frac{1}{4!} \left(x-\frac\pi2\right)^4
    - \frac{1}{6!} \left(x-\frac\pi2\right)^6
    $.
    \end{red}
    \item How is your result in (a) similar to the degree 6 Taylor polynomial of $\cos(x)$ that is centered at $a=0$?
    
    \begin{red}
    Well hey, it's totally the same, except all of the $x$'s have been replaced by $\left(x-\frac{\pi}{2}\right)$'s.
    \end{red}
    \item Recall the trigonometric identity that states $\sin(x) = \cos(x-\tfrac\pi2)$. How does this identity help you explain what you found in (a) and (b)?
    
    \begin{red}
    This is actually a perfect description of what has happened: we took the thing we got for $\cos(x)$ and just replaced $x$ by $(x-\tfrac\pi2)$, and according to this trigonometric identity, that gives us the values of $\sin(x)$.
    \end{red}
    \end{enumerate}

\item 

\begin{enumerate}
    \item (\href{https://faculty.gvsu.edu/boelkinm/Home/ACS/sec-8a-2-taylor-poly.html#ez-8a-2-3}{AC 8a.2 Exercise 8}) Note that $f(0.5) = \ln(1.5)$. Use $T_5(x)$ to estimate $\ln(1.5)$.
    
    \begin{red} $\ln(1.5) \approx T_5(0.5) = .40729166666\ldots$.
    \end{red}
    
    \item Observe that $g(1.5) = \ln(1.5)$. Use $P_5(x)$ to estimate $\ln(1.5)$.
    
    \begin{red}$\ln(1.5) \approx P_5(1.5) = .40729166666\ldots$.
    \end{red}
    
    \item Are the estimates of $\ln(1.5)$ generated by $T_5(x)$ and $P_5(x)$ the same or different? Why do you think this happens?
    
    \begin{red}
    They're the same! Since $g(x)$ and $f(x)$ are just horizontal shifts of each other, their Taylor polynomials should agree as long as we compensate for the horizontal shift by recentering the approximation.
    \end{red}
\end{enumerate}

\end{enumerate}
\end{document}